\documentclass[reprint,amsmath,amssymb,aps,floatfix]{revtex4-2}
\usepackage{graphicx}
\usepackage{dcolumn}
\usepackage{bm}
\usepackage{amsmath}
\usepackage{amssymb}
\usepackage{hyperref}
\begin{document}
\preprint{APS/Academic Review}
\title{Topic: terahertz generation -- 文献综述 v4.0}
\author{Claude Code Academic Brain}
\affiliation{%
 Self Learning Project%
}%
\date{2026-04-28}%
\begin{abstract}
本综述基于OpenAlex数据库检索的30篇论文，覆盖6个主题方向。
通过Thematic Synthesis方法，对各技术路线的核心权衡与研究空白进行系统分析。
\end{abstract}
\maketitle
\section{引言}
\par\par{太赫兹(THz)波段（0.1-10 THz）位于微波与红外之间，是连接电子学与光学的桥梁，在成像、光谱学、通信和安全检测等领域具有重要应用前景。然而，THz辐射的高效产生一直是该领域核心技术难题，制约着相关应用的发展。}
\par\par{【研究空白】现有文献存在以下不足：}
\begin{itemize}
\item PCA材料：缺乏载流子寿命<100fs且迁移率>5000 cm²/Vs的材料方案
\item PCA结构：缺乏对大面积天线与超宽带之间物理极限的系统理论研究
\item 光整流晶体：缺乏同时实现>5 THz带宽 AND >100μJ能量输出的晶体方案
\item 空气等离子体：缺乏对远程传输中相位稳定性问题的系统性解决方案
\end{itemize}
\par\par{本综述基于OpenAlex数据库检索的30篇相关论文，对7个主题方向进行系统性的主题综合，识别各技术路线的核心权衡与研究空白，为不同应用场景下的THz源选择提供参考依据。}
\section{方法与结果}
\subsection{PCA材料}
\par\par{PCA材料领域的核心科学问题在于：何种材料能在载流子寿命、迁移率、暗电阻三方面同时满足PCA需求？}

\par\par{当前研究主要采用InAs、GaAs等技术路线。}
\par\par{Michael B. Johnston、D. M.等人(2002)与Nezih Tolga Yardimci、Shan等人(2015)的工作对比表明：}
虽都聚焦于InAs方向，但Michael B. Johnston、D. M.等人侧重于Simulation of terahertz generation，而Nezih Tolga Yardimci、Shan等人则聚焦于High-Power Terahertz Generation Usi，在具体研究对象上存在差异。
\par\par{【共识】LT-GaAs仍是超宽带PCA的主流选择；载流子寿命<1ps是实现高带宽的共识}

\par\par{【研究空白】该领域存在以下未被充分研究的问题：}
\begin{itemize}
\item 缺乏载流子寿命<100fs且迁移率>5000 cm²/Vs的材料方案
\item 缺乏对新型纳米复合材料的长期稳定性的系统研究
\item 近年该主题论文数量有所下降，可能存在新的技术突破点
\end{itemize}

\subsubsection*{性能对比}
\begin{table}[h!]
\caption{关键工作对比: PCA材料}
\begin{tabular}{@{}llp{2.5cm}p{2.5cm}r@{}}
\hline
Author(s) & Year & Method & Key Metrics & Citations\\
\hline
Michael B. Johnston、 & 2002 & InAs & --- & 360\\
Nezih Tolga Yardimci & 2015 & --- & 240.0MW、3.8MW & 217\\
G. Kh. Kitaeva & 2008 & --- & 10.0THz & 180\\
\hline
\end{tabular}
\end{table}

\subsubsection*{未来方向}
基于上述研究空白，未来研究应重点关注：缺乏载流子寿命<100fs且迁移率>5000 cm²/Vs的材料方案

\subsection{PCA结构}
\par\par{PCA结构领域的核心科学问题在于：天线几何结构如何影响阻抗匹配、带宽和光泵吸收效率的权衡？}

\par\par{当前研究主要采用dipole、large-area、plasmonic等技术路线。}
\par\par{Michael B. Johnston、D. M.等人(2002)与Santosh Jain、J. Parashar等人(2013)的工作对比表明：}
虽都聚焦于dipole方向，但Michael B. Johnston、D. M.等人侧重于Simulation of terahertz generation，而Santosh Jain、J. Parashar等人则聚焦于Effect of magnetic field on，在具体研究对象上存在差异。
\par\par{【共识】大面积 plasmonic 电极可显著提升泵浦效率；bow-tie 是宽带应用的主流结构}

\par\par{【研究空白】该领域存在以下未被充分研究的问题：}
\begin{itemize}
\item 缺乏对大面积天线与超宽带之间物理极限的系统理论研究
\item 缺乏在真实工作条件下(高泵浦功率密度)的热管理对比研究
\item 近年该主题论文数量有所下降，可能存在新的技术突破点
\end{itemize}

\subsubsection*{性能对比}
\begin{table}[h!]
\caption{关键工作对比: PCA结构}
\begin{tabular}{@{}llp{2.5cm}p{2.5cm}r@{}}
\hline
Author(s) & Year & Method & Key Metrics & Citations\\
\hline
Michael B. Johnston、 & 2002 & dipole & --- & 360\\
Santosh Jain、J. Para & 2013 & dipole & --- & 222\\
Nezih Tolga Yardimci & 2015 & dipole & 240.0MW、3.8MW & 217\\
\hline
\end{tabular}
\end{table}

\subsubsection*{未来方向}
基于上述研究空白，未来研究应重点关注：缺乏对大面积天线与超宽带之间物理极限的系统理论研究

\subsection{光整流晶体}
\par\par{光整流晶体领域的核心科学问题在于：相位匹配与损伤阈值的矛盾如何限制THz输出能量与带宽的提升？}

\par\par{当前研究主要采用ZnTe、tilted pulse front、LiNbO3、GaSe等技术路线。}
\par\par{Anastasios D. Koulouklidi等人(2020)与Matthias C. Hoffmann、Ka-L等人(2007)的工作对比表明：}
Anastasios D. Koulouklidi等人(2020)实现了2.36%，而Matthias C. Hoffmann、Ka-L等人(2007)则达到2.5THz，两者在指标上相差0.1%。
\par\par{【共识】倾斜脉冲前技术是LiNbO3中实现相位匹配的主流方案；有机晶体在带宽上具有优势}

\par\par{【研究空白】该领域存在以下未被充分研究的问题：}
\begin{itemize}
\item 缺乏同时实现>5 THz带宽 AND >100μJ能量输出的晶体方案
\item 有机晶体在强泵浦下的光稳定性数据几乎空白
\item 缺乏对晶体热效应的系统性理论和实验研究
\end{itemize}

\subsubsection*{性能对比}
\begin{table}[h!]
\caption{关键工作对比: 光整流晶体}
\begin{tabular}{@{}llp{2.5cm}p{2.5cm}r@{}}
\hline
Author(s) & Year & Method & Key Metrics & Citations\\
\hline
Anastasios D. Koulou & 2020 & --- & 2.36%、3.9μm & 296\\
Matthias C. Hoffmann & 2007 & LiNbO3 & 2.5THz、400.0μJ & 194\\
G. Kh. Kitaeva & 2008 & GaSe & 10.0THz & 180\\
\hline
\end{tabular}
\end{table}

\subsubsection*{未来方向}
基于上述研究空白，未来研究应重点关注：缺乏同时实现>5 THz带宽 AND >100μJ能量输出的晶体方案

\subsection{空气等离子体}
\par\par{空气等离子体领域的核心科学问题在于：空气等离子体方法如何在远程、高峰值功率、高转换效率三方面实现协同提升？}

\par\par{当前研究主要采用filamentation、terawatt、two-color等技术路线。}
\par\par{Anastasios D. Koulouklidi等人(2020)与Indranuj Dey、Kamalesh Jan等人(2017)的工作对比表明：}
Anastasios D. Koulouklidi等人(2020)实现了2.36%，而Indranuj Dey、Kamalesh Jan等人(2017)则达到50.0THz，两者在指标上相差47.6%。
虽都聚焦于filamentation方向，但Anastasios D. Koulouklidi等人侧重于Observation of extremely efficient，而Indranuj Dey、Kamalesh Jan等人则聚焦于Highly efficient broadband terahert，在具体研究对象上存在差异。
\par\par{【共识】双光束方案是实现高效率的主流选择；空气等离子体是远程THz产生的唯一可行方案}

\par\par{【研究空白】该领域存在以下未被充分研究的问题：}
\begin{itemize}
\item 缺乏对远程传输中相位稳定性问题的系统性解决方案
\item 缺乏在真实大气条件下(湿度、湍流)THz传输的系统性数据
\item 缺乏高重频(>1kHz)空气等离子体THz源的可行性研究
\end{itemize}

\subsubsection*{性能对比}
\begin{table}[h!]
\caption{关键工作对比: 空气等离子体}
\begin{tabular}{@{}llp{2.5cm}p{2.5cm}r@{}}
\hline
Author(s) & Year & Method & Key Metrics & Citations\\
\hline
Anastasios D. Koulou & 2020 & filamentation & 2.36%、3.9μm & 296\\
Indranuj Dey、Kamales & 2017 & filamentation & 50.0THz & 183\\
Taek Il Oh、Yong Sing & 2013 & filamentation & 60.0THz、0.1mJ & 179\\
\hline
\end{tabular}
\end{table}

\subsubsection*{未来方向}
基于上述研究空白，未来研究应重点关注：缺乏对远程传输中相位稳定性问题的系统性解决方案

\subsection{QCL激光器}
\par\par{QCL激光器领域的核心科学问题在于：QCL如何在突破4 THz频段限制、降低冷却要求、提升输出功率三个方向上取得突破？}

\par\par{当前研究主要采用heterostructure、mid-infrared、quantum cascade等技术路线。}
\par\par{Karun Vijayraghavan、Yifan等人(2013)与Chao Zhou、Y. P. Liu等人(2018)的工作对比表明：}
在技术路径上，quantum cascade与heterostructure代表了两条不同的方向，分别体现了高频段(>4 THz)与室温工作的权衡。
\par\par{【共识】QCL是唯一实现电泵浦的固态THz源；向室温工作是明确的发展方向}

\par\par{【研究空白】该领域存在以下未被充分研究的问题：}
\begin{itemize}
\item 缺乏4 THz以上、室温连续波工作的可行性理论分析
\item 缺乏对QCL长期工作稳定性的系统性研究
\item 缺乏对QCL与外部高效耦合器的宽带匹配研究
\end{itemize}

\subsubsection*{性能对比}
\begin{table}[h!]
\caption{关键工作对比: QCL激光器}
\begin{tabular}{@{}llp{2.5cm}p{2.5cm}r@{}}
\hline
Author(s) & Year & Method & Key Metrics & Citations\\
\hline
Karun Vijayraghavan、 & 2013 & quantum cascade & --- & 193\\
Chao Zhou、Y. P. Liu & 2018 & heterostructure & 5.0THz & 185\\
Leonhard Prechtel、Li & 2012 & --- & 1.0THz、640.0GHz & 171\\
\hline
\end{tabular}
\end{table}

\subsubsection*{未来方向}
基于上述研究空白，未来研究应重点关注：缺乏4 THz以上、室温连续波工作的可行性理论分析

\subsection{检测技术}
\par\par{检测技术领域的核心科学问题在于：不同检测方法在灵敏度、带宽、动态范围之间如何权衡？}

\par\par{当前研究主要采用coherent detection、supercontinuum generation、broadband generation、metamaterial等技术路线。}
\par\par{Ki‐Yong Kim、Antoinette J.等人(2008)与Liang Luo、Ioannis Chatzak等人(2014)的工作对比表明：}
在技术路径上，coherent detection与metamaterial代表了两条不同的方向，分别体现了相干探测(电光采样)与能量探测(热敏)的权衡。
\par\par{【共识】电光采样是宽带相干探测的主流；低温bolometer是最高灵敏度方案}

\par\par{【研究空白】该领域存在以下未被充分研究的问题：}
\begin{itemize}
\item 缺乏对超宽带(>20 THz)实时探测的系统性技术方案
\item 缺乏对室温下同时实现高灵敏度( NEP<10 pW/√Hz)和宽带(>5 THz)的检测器的可行性分析
\item 缺乏对大规模阵列探测器的系统性研究
\end{itemize}

\subsubsection*{性能对比}
\begin{table}[h!]
\caption{关键工作对比: 检测技术}
\begin{tabular}{@{}llp{2.5cm}p{2.5cm}r@{}}
\hline
Author(s) & Year & Method & Key Metrics & Citations\\
\hline
Ki‐Yong Kim、Antoinet & 2008 & coherent detection & --- & 822\\
Liang Luo、Ioannis Ch & 2014 & metamaterial & --- & 214\\
F. Brunner、O‐Pil Kwo & 2008 & --- & 2.5THz、2.5THz & 172\\
\hline
\end{tabular}
\end{table}

\subsubsection*{未来方向}
基于上述研究空白，未来研究应重点关注：缺乏对超宽带(>20 THz)实时探测的系统性技术方案
\section{讨论}
\par\par{本综述对THz产生的多条技术路线进行了系统性的主题综合分析，揭示了各方法在效率、带宽、功率三个核心指标上的权衡关系。}
\par\par{光电导天线在超宽带应用场景具有独特优势，但受限于材料载流子动力学；光整流在LiNbO3倾斜脉冲前技术支撑下实现了mJ级THz输出，但有机晶体的高带宽优势尚未转化为实际能量输出；空气等离子体方法突破了远程产生的物理限制，但相位稳定性仍是实用化的瓶颈；QCL正在向室温高功率方向发展，但4 THz以上频段仍是挑战；超表面为THz器件的平面化集成提供了新途径，但距离高效THz源仍有差距。}
\par\par{从研究空白来看，缺乏对各技术路线在不同应用场景下的系统性性能边界分析，这将是未来研究的重要方向。}
\appendix
\section{附录：论文列表}
\begin{table}[h!]
\caption{检索到的论文列表}
\label{tab:papers}
\begin{tabular}{@{ }p{10cm}r@{ }}
\hline
\# & Title & Citations\\
\hline
1 & Simulation of terahertz generation at semiconductor surfaces & 360\\
2 & Wavelength Scaling of Terahertz Generation by Gas Ionization & 382\\
3 & Broadband terahertz generation from metamaterials & 214\\
4 & Terahertz generation by means of optical lasers & 180\\
5 & Broadly tunable terahertz generation in mid-infrared quantum cascade lasers & 193\\
6 & Observation of extremely efficient terahertz generation from mid-infrared two-co... & 296\\
7 & Coherent control of terahertz supercontinuum generation in ultrafast laser–gas i... & 822\\
8 & Efficient terahertz generation by optical rectification at 1035 nm & 194\\
9 & Highly efficient broadband terahertz generation from ultrashort laser filamentat... & 183\\
10 & Broadband Terahertz Generation via the Interface Inverse Rashba-Edelstein Effect & 185\\
11 & Terahertz generation using plasmonic photoconductive gratings & 147\\
12 & Effect of magnetic field on terahertz generation via laser interaction with a ca... & 222\\
13 & High-Power Terahertz Generation Using Large-Area Plasmonic Photoconductive Emitt... & 217\\
14 & Spectral phase control of interfering chirped pulses for high-energy narrowband ... & 152\\
15 & Time-resolved ultrafast photocurrents and terahertz generation in freely suspend... & 171\\
16 & Intense terahertz generation in two-color laser filamentation: energy scaling wi... & 179\\
17 & A hydrogen-bonded organic nonlinear optical crystal for high-efficiency terahert... & 172\\
18 & Experimental Comparison of UTC- and PIN-Photodiodes for Continuous-Wave Terahert... & 89\\
19 & Data Mining for Terahertz Generation Crystals & 56\\
20 & Continuous Wave Terahertz Generation From Ultra-Fast InP-Based Photodiodes & 92\\
21 & Pulse sequences for efficient multi-cycle terahertz generation in periodically p... & 92\\
22 & Toward remote high energy terahertz generation & 74\\
23 & Monolithic dual-mode distributed feedback semiconductor laser for tunable contin... & 129\\
24 & Sub‐terahertz and terahertz generation in long‐wavelength quantum cascade lasers & 65\\
25 & Boosting Terahertz Generation in Laser-Field Ionized Gases Using a Sawtooth Wave... & 111\\
26 & Terahertz generation from graphite & 52\\
27 & Terahertz Generation by Dynamical Photon Drag Effect in Graphene Excited by Femt... & 110\\
28 & Pulse sequences for efficient multi-cycle terahertz generation in periodically p... & 90\\
29 & Narrowband terahertz generation with chirped-and-delayed laser pulses in periodi... & 80\\
30 & Terahertz generation from laser-induced plasma & 62\\
\hline
\end{tabular}
\end{table}
\end{document}